\chapter{Echelon Form}\label{ch:echelon}

Chapter~\ref{ch:systems} proved that elimination does not change the answer. This chapter proves
that it arrives at one: the output is in reduced row echelon form, from which the solution set can
be read off. On paper the proof is ``keep going until it is''; in Lean it is an invariant carried
column by column, and the one lemma that made the invariant provable is about a detail the paper
proof never mentions.

\section{The definition}

\begin{definition}[Reduced row echelon form]
  A matrix is in reduced row echelon form (RREF) if:
  \begin{enumerate}[nosep]
    \item every zero row is below every nonzero row;
    \item the first nonzero entry of each nonzero row (its \emph{pivot}) is $1$;
    \item the pivot of each nonzero row is strictly to the right of the pivot of the row above;
    \item each pivot is the only nonzero entry in its column.
  \end{enumerate}
\end{definition}

Conditions 1--3 alone define \emph{row echelon form}; condition 4 is the ``reduced''. From an
RREF matrix the solutions are immediate: the pivot columns' variables are determined by the
others (the \emph{free} variables), one equation each, and the free variables are arbitrary. The
number of pivots is the rank.

\begin{theorem}[Existence]
  Every matrix can be brought to RREF by elementary row operations.
\end{theorem}
\begin{proof}
  Induction on the columns, maintaining: the first $p$ rows are in RREF with their pivots in
  the columns seen so far, and every later row is zero in those columns. Initially $p = 0$. For
  the next column $c$: if rows $p, p+1, \ldots$ are all zero in column $c$, the invariant holds
  with $c$ added and $p$ unchanged. Otherwise pick a row $q \ge p$ with a nonzero entry, swap it
  to row $p$, scale row $p$ to make the entry $1$, and add multiples of row $p$ to every other row
  to clear column $c$. Rows above $p$ keep their pivots (their pivot columns are left of $c$,
  where row $p$ is zero); rows below $p$ are now zero through column $c$; row $p$ is a pivot row
  with pivot at $c$. So the invariant holds with $p + 1$.
\end{proof}

\begin{theorem}[Uniqueness]
  The RREF of a matrix is unique.
\end{theorem}

Uniqueness is true and not proved by the engine; the standard argument is that two RREF matrices
with the same row space have the same pivot columns and then the same rows. The engine's claim is
existence, for its own algorithm.

\section{The algorithm}

\begin{minted}{lean}
def columnBody (m : Mat) (col p q : Nat) : List Op :=
  let swaps := if q = p then [] else [Op.swap q p]
  let m1 := runOps swaps m
  let piv := entry m1 p col
  let scales := if piv = 1 then [] else [Op.scale p piv.inv]
  let m2 := runOps scales m1
  let factors := ((List.range m.length).filter fun i => i ≠ p ∧ entry m2 i col ≠ 0).map fun i => (i, -(entry m2 i col))
  swaps ++ scales ++ factors.map fun x => Op.addMul x.1 p x.2

def columnOps (m : Mat) (col p : Nat) : Option (List Op) :=
  ((List.range m.length).find? fun i => p ≤ i && entry m i col != 0).map (columnBody m col p)

def go (m : Mat) (col p : Nat) : Nat → List (Nat × Op)
  | 0 => []
  | k + 1 =>
    if p ≥ m.length then [] else
    match columnOps m col p with
    | none => go m (col + 1) p k
    | some ops => let tagged := ops.map (col, ·); tagged ++ go (run tagged m) (col + 1) (p + 1) k
\end{minted}

This is the proof's induction step as code: swap, scale, clear; or move on. The recursion is on
the number of columns remaining, so it is structural and needs no fuel.

\section{The invariant}

\begin{minted}{lean}
structure Inv (m : Mat) (p col : Nat) (pivs : List Nat) : Prop where
  len    : pivs.length = p
  lt     : ∀ c ∈ pivs, c < col
  sorted : pivs.Pairwise (· < ·)
  pivot  : ∀ i, i < p → entry m i (pivs.getD i 0) = 1
  lead   : ∀ i, i < p → ∀ j, j < pivs.getD i 0 → entry m i j = 0
  alone  : ∀ i, i < p → ∀ k, k ≠ i → entry m k (pivs.getD i 0) = 0
  below  : ∀ i, p ≤ i → ∀ j, j < col → entry m i j = 0
  p_le   : p ≤ m.length

def IsRref (m : Mat) (w : Nat) : Prop := ∃ p pivs, Inv m p w pivs
theorem rref_isRref (m : Mat) : IsRref (rref m) (ncols m)
\end{minted}

Match the clauses to the definition: \lc{pivot} and \lc{lead} are condition 2 (a leading $1$);
\lc{sorted} is condition 3; \lc{alone} is condition 4; \lc{below} is condition 1 in the form
``rows $\ge p$ are zero in every column seen''. When \lc{col} reaches the width, \lc{below} says
the later rows are entirely zero, and \lc{IsRref} is the definition.

The four transition lemmas are the four sentences of the paper proof: \lc{Inv.next_none} (no
pivot: the column moves on), \lc{Inv.swap_ge} (swapping row $p$ with a row below keeps the
invariant, since both are zero left of \lc{col}), \lc{Inv.scale_p} (scaling row $p$ keeps it),
and \lc{Inv.next_some} (after clearing, one more pivot row). \lc{go_inv} folds them over the
recursion.

\section{The lemma the paper proof skips}

The clearing step applies one row addition per nonzero entry, \emph{sequentially}. The paper proof
says ``clear the column'' as if all the additions happened at once, and for the invariant that
is what one wants: every row other than $p$ changes by its own multiple of row $p$, and row $p$
itself does not change. Sequential application agrees with simultaneous application only because
each addition touches a different row and reads row $p$, which none of them touches. The lemma
that says so is:

\begin{minted}{lean}
theorem entry_clears {p : Nat} : ∀ (L : List (Nat × Rat)) (m : Mat), p < m.length →
    (L.map Prod.fst).Nodup → (∀ x ∈ L, x.1 ≠ p) → (∀ x ∈ L, x.1 < m.length) →
    ∀ k l, entry (runOps (L.map fun x => Op.addMul x.1 p x.2) m) k l =
      match coeffOf k L with
      | some c => entry m k l + c * entry m p l
      | none => entry m k l
\end{minted}

The factors are read from the matrix \emph{before} any clearing (\lc{factors} is computed from
\lc{m2}, not updated as the loop runs), and the list of rows has no duplicates and excludes $p$.
With this lemma every clause of \lc{Inv} after the column step is checked one entry at a time,
and the proof never follows the loop. It is the kind of lemma one only discovers by trying to
prove the invariant without it.

\section{Two remarks on the statement}

The theorem is stated in the width of the \emph{input}, \lc{ncols m}, rather than of the output.
That avoids proving that row operations preserve rectangularity, which the padding convention of
Chapter~\ref{ch:systems} makes false in general (a padded row may be longer). For the rectangular
matrices the parser admits, the two widths coincide.

And the field algebra is again \lc{grind}'s. \lc{ring} is Mathlib's tactic, unavailable in the
Init-only engine, and the rational identities that clearing needs --- that $a + (-a/b)\cdot b = 0$
and the like --- are within \lc{grind}'s reach. The whole proof is 412 lines and imports nothing.

\begin{logbox}[Before the proof]
  The first version decided at run time whether the output was in RREF, and refused the
  evaluation if not. That is the honest boundary of Chapter~\ref{ch:order}: check what the
  proof needs, fail loudly. It was also a promise to come back, and the check is gone now.
\end{logbox}

\begin{trybox}
\begin{minted}{text}
rref([0,0,1;1,0,0])
rref([2,4,6;1,2,3])
rref([1,1,1;1,2,3;1,4,9])
\end{minted}
  The first needs a swap. The second has rank one. The third is a Vandermonde matrix; its RREF
  is the identity, which is why polynomial interpolation through three points has one answer.
\end{trybox}

\section{Exercises}

\begin{exercisebox}[Reading off solutions]
  Bring $\begin{pmatrix} 1 & 2 & 0 & 3 \\ 2 & 4 & 1 & 8 \end{pmatrix}$ to RREF by hand, name the
  pivot and free variables, and write the solution set as a point plus a span.
\end{exercisebox}

\begin{exercisebox}[Uniqueness]
  Prove that two matrices in RREF with the same row space are equal. (First show they have the
  same pivot columns, by considering the leftmost column in which they differ.)
\end{exercisebox}

\begin{exercisebox}[Why the factors are read first]
  Construct a $3 \times 1$ matrix on which reading the clearing factors \emph{after} each
  addition, rather than before, would clear the column incorrectly if the pivot row were also
  in the list.
\end{exercisebox}
